\section{The shared-decoder model}\label{sec:model}
Let $d\ge1$ and write $I=\Z/d\Z$. For an entry alphabet $A$, the complete
state space is $\mathcal X=A^{I\times I}$. In the linear setting $A=F^b$,
where $F$ is a field and $b\ge1$; thus $N=bd^2$ is the dimension over $F$.
For binary tape inputs the bit length is $L=d^2b$. The canonical width in
\textup{(T)} is $b(d)=\max(1,\lceil\log_2d\rceil)$, but algebraic statements
will quantify over every positive $b$.

Define the column and row shifts and the transpose by
\[
 (T^sX)_{i,j}=X_{i,j+s},\qquad
 (R^sX)_{i,j}=X_{i+s,j},\qquad
 (\tau X)_{i,j}=X_{j,i}.
\]
On a row or column word, $\sigma^s$ shifts its single cyclic index in the same
orientation. Thus $\tau T^s=R^s\tau$, $\tau R^s=T^s\tau$, and $T^sR^t=R^tT^s$.

\begin{definition}[A shared rule and its correctness source]\label{def:shared-rule}
\formal{Model.lean}{ClassUniform,IsSharedCode}
A receiver partition $p:I\to\{1,\ldots,M\}$ uses every label.
A complete shared rule $K:\mathcal X\to\mathcal X$ has total functions
$f_J:\mathcal X\to A^I$ such that
\[
 \operatorname{col}_c K(X)=f_{p(c)}(T^cX).
\]
The decoder has no additional receiver-index argument. In the linear setting,
a subspace $C\le\mathcal X$ is a correctness source if
$f_{p(c)}(T^cX)=\operatorname{row}_cX$ for every $X\in C$ and every $c$.
Its dimension is denoted $D$.
\end{definition}

The predicates are distinct: a correctness source need not contain arbitrary
intermediate states of a composition. A linear source does not imply that its
decoders are linear. A \emph{linear complete rule} means that $K$ itself is
$F$-linear on all of $\mathcal X$. All definitions permit mixing between lanes.
A map $H$ is column-equivariant if $HT^s=T^sH$ for every $s$; one-class
complete rules are exactly these maps. Equivariance does not assert that the
map is cheap to compute on tapes.

All complete factorizations in Part II act on this same state space and use
this same column-view action, unless a theorem explicitly says otherwise.
For a finite alphabet, exact factorization of the bijection $\tau$ by total
square maps forces every factor to be bijective. Correctness on a proper
source, or a larger workspace with a promised initialization, does not imply
this conclusion. Algebraic row or column operations incur no native time
charge until a separately stated implementation supplies it.
